\documentclass[11pt]{article}

\usepackage[T1]{fontenc}
\usepackage{lmodern}
\usepackage{microtype}
\usepackage[a4paper,margin=30mm]{geometry}
\usepackage{amsmath,amssymb,amsthm,mathtools}
\usepackage{enumitem}
\usepackage[hidelinks]{hyperref}
\usepackage{bookmark}
\hypersetup{
  pdftitle={Rational and Integral Points on y2 - x2 y + z2 + 1 = 0},
  pdfauthor={Jamal Agbanwa},
  pdfsubject={A complete self-contained classification of the rational solutions and a proof that no integral solutions exist}
}

\allowdisplaybreaks
\numberwithin{equation}{section}

\newtheorem{theorem}{Theorem}[section]
\newtheorem{proposition}[theorem]{Proposition}
\newtheorem{lemma}[theorem]{Lemma}
\newtheorem{corollary}[theorem]{Corollary}
\theoremstyle{definition}
\newtheorem{definition}[theorem]{Definition}
\newtheorem{example}[theorem]{Example}
\theoremstyle{remark}
\newtheorem{remark}[theorem]{Remark}

\newcommand{\Q}{\mathbb{Q}}
\newcommand{\Z}{\mathbb{Z}}
\newcommand{\Pone}{\mathbb{P}^{1}}
\newcommand{\ord}{\operatorname{ord}}
\newcommand{\eps}{\varepsilon}

\title{Rational and Integral Points on\\[3pt]
\texorpdfstring{$y^{2}-x^{2}y+z^{2}+1=0$}{y^2 - x^2 y + z^2 + 1 = 0}}
\author{Jamal Agbanwa}
\date{18 July 2026}

\begin{document}
\maketitle

\begin{abstract}
We give a complete, unconditional, and self-contained description of the
rational solutions of
\[
 y^{2}-x^{2}y+z^{2}+1=0.
\]
If a nonzero rational number is written in lowest terms as
$x=\eps a/b$, where $a,b>0$, $\gcd(a,b)=1$, and $\eps\in\{\pm1\}$,
then the fibre above $x$ has a rational point if and only if
$N_{a,b}=a^{4}-4b^{4}$ is positive and every prime congruent to $3$
modulo $4$ occurs in $N_{a,b}$ with even exponent.  Whenever this
criterion holds, an integral representation $N_{a,b}=C^{2}+D^{2}$
exists, and a homogeneous rational parametrisation of the entire fibre
is written down explicitly.  We then prove that every admissible
reduced fraction $a/b$ has odd numerator and even denominator.  It
follows, more strongly than the absence of integral triples, that no
rational solution can have integral $x$.  A second independent
congruence proof of the integral obstruction is included.  Finally, we
exhibit both a complete one-parameter fibre over $x=\pm3/2$ and a Pell
recurrence producing infinitely many distinct rational
$x$-coordinates.
\end{abstract}

\section{Statement of the classification}

Let
\[
 \mathcal S(F)=\bigl\{(x,y,z)\in F^{3}:
 y^{2}-x^{2}y+z^{2}+1=0\bigr\}
\]
for a subfield $F\subseteq\mathbb R$.  For coprime positive integers
$a,b$, put
\[
 N_{a,b}=a^{4}-4b^{4}.
\]
For a prime $p$ and a nonzero integer $n$, let $\ord_{p}(n)$ denote the
exponent of $p$ in the prime factorisation of $|n|$.

\begin{theorem}[Complete rational classification]\label{thm:main}
Let $a,b\in\Z_{>0}$ satisfy $\gcd(a,b)=1$, and let
$\eps\in\{\pm1\}$.  The fibre
\[
 \mathcal S_{\eps a/b}(\Q)
 =\{(y,z)\in\Q^{2}:(\eps a/b,y,z)\in\mathcal S(\Q)\}
\]
is nonempty if and only if
\begin{equation}\label{eq:criterion}
 N_{a,b}>0
 \quad\text{and}\quad
 \ord_{q}(N_{a,b})\equiv0\pmod 2
 \quad\text{for every prime }q\equiv3\pmod4.
\end{equation}

Assume \eqref{eq:criterion}.  Choose integers $C,D$ such that
\begin{equation}\label{eq:CD}
 C^{2}+D^{2}=N_{a,b}.
\end{equation}
For $[u:v]\in\Pone(\Q)$, set
\[
 U=u^{2}+v^{2}.
\]
Then $U\ne0$, and every rational point in the fibre, exactly once, is
obtained from
\begin{align}
 x&=\eps\frac ab,\label{eq:param-x}\\
 y&=\frac{a^{2}U+C(u^{2}-v^{2})-2Duv}
 {2b^{2}U},\label{eq:param-y}\\
 z&=\frac{D(u^{2}-v^{2})+2Cuv}
 {2b^{2}U}.\label{eq:param-z}
\end{align}
Here ``exactly once'' refers to a fixed choice of $a,b,\eps,C,D$ and
the usual projective identification $[u:v]=[\lambda u:\lambda v]$ for
$\lambda\in\Q^{\times}$.

Moreover, every pair $(a,b)$ satisfying \eqref{eq:criterion} obeys
\begin{equation}\label{eq:parity-conclusion}
 a\equiv1\pmod2,
 \qquad
 b\equiv0\pmod2.
\end{equation}
Consequently no point of $\mathcal S(\Q)$ has integral $x$, and in
particular
\[
 \mathcal S(\Z)=\varnothing.
\]
\end{theorem}

The proof is assembled from elementary identities, the two-square
criterion, and the rational parametrisation of a circle.  All of these
ingredients are proved below.

\section{Elementary reductions}

\begin{proposition}[Completion of the square]\label{prop:completion}
For all real $x,y,z$ one has the identity
\begin{equation}\label{eq:master-identity}
4\bigl(y^{2}-x^{2}y+z^{2}+1\bigr)
=(2y-x^{2})^{2}+(2z)^{2}-(x^{4}-4).
\end{equation}
Thus
\begin{equation}\label{eq:norm-equation}
 y^{2}-x^{2}y+z^{2}+1=0
 \quad\Longleftrightarrow\quad
 (2y-x^{2})^{2}+(2z)^{2}=x^{4}-4.
\end{equation}
Every rational solution satisfies
\begin{equation}\label{eq:real-inequalities}
 0<y<x^{2}
 \qquad\text{and}\qquad
 x^{4}>4.
\end{equation}
\end{proposition}

\begin{proof}
Expanding the right-hand side of \eqref{eq:master-identity} gives
\[
4y^{2}-4x^{2}y+x^{4}+4z^{2}-x^{4}+4,
\]
which is the left-hand side.  This proves \eqref{eq:norm-equation}.

If the original equation holds, then
\begin{equation}\label{eq:product-form}
 y(x^{2}-y)=z^{2}+1>0.
\end{equation}
If $y\le0$, then $x^{2}-y\ge0$, so the left-hand side of
\eqref{eq:product-form} is nonpositive, a contradiction.  Hence $y>0$;
then \eqref{eq:product-form} forces $x^{2}-y>0$.  This proves
$0<y<x^{2}$.

Equation \eqref{eq:norm-equation} gives $x^{4}\ge4$.  Equality would
force $2y-x^{2}=0$ and $z=0$, and hence $x^{2}=2$.  A rational number
cannot have square $2$: if $x=r/s$ in lowest terms and $r^{2}=2s^{2}$,
then $r$ is even, after which $s$ is also even, contradicting
$\gcd(r,s)=1$.  Therefore a rational solution has $x^{4}>4$.
\end{proof}

\section{The two-square criterion over \texorpdfstring{$\Q$}{Q}}

We prove the precise arithmetic criterion needed later.  The argument
also proves that allowing rational squares does not enlarge the class
of positive integers representable as a sum of two squares.

\begin{lemma}\label{lem:3mod4-divides-both}
Let $q\equiv3\pmod4$ be prime.  If $A,B\in\Z$ and
$q\mid A^{2}+B^{2}$, then $q\mid A$ and $q\mid B$.
\end{lemma}

\begin{proof}
If $q\mid B$, then $q\mid A^{2}$ and hence $q\mid A$.  Suppose instead
that $q\nmid B$.  Modulo $q$, let $r\equiv AB^{-1}$.  Then
$r^{2}\equiv-1\pmod q$.

For completeness, multiplication by the nonzero residue $r$ permutes
the nonzero residues modulo $q$.  Multiplying all these residues and
cancelling the nonzero product gives
$r^{q-1}\equiv1\pmod q$.  On the other hand,
\[
 r^{q-1}=(r^{2})^{(q-1)/2}
 \equiv(-1)^{(q-1)/2}\equiv-1\pmod q,
\]
because $(q-1)/2$ is odd.  This contradiction proves that $q\mid B$,
and then also $q\mid A$.
\end{proof}

\begin{corollary}\label{cor:valuation-even}
If $q\equiv3\pmod4$ is prime and $A,B\in\Z$ are not both zero, then
\[
 \ord_{q}(A^{2}+B^{2})
\]
is even.
\end{corollary}

\begin{proof}
Let
\[
 e=\min\{\ord_{q}(A),\ord_{q}(B)\},
\]
where $\ord_q(0)$ may be regarded as $+\infty$.  Write
$A=q^{e}A_{0}$ and $B=q^{e}B_{0}$, with at least one of $A_{0},B_{0}$
not divisible by $q$.  By Lemma~\ref{lem:3mod4-divides-both},
$q\nmid A_{0}^{2}+B_{0}^{2}$.  Therefore
\[
 \ord_{q}(A^{2}+B^{2})=2e.
\]
\end{proof}

\begin{lemma}[Primes congruent to $1$ modulo $4$]\label{lem:fermat-two-square}
Every prime $p\equiv1\pmod4$ can be written as a sum of two integer
squares.
\end{lemma}

\begin{proof}
Write $p=4k+1$ and put
\[
 m=\frac{p-1}{2}=2k.
\]
We first produce a square root of $-1$ modulo $p$.  Wilson's congruence
\begin{equation}\label{eq:wilson}
 (p-1)!\equiv-1\pmod p
\end{equation}
follows by pairing every nonzero residue with its multiplicative
inverse: the only self-inverse nonzero residues are $1$ and $-1$, since
$t^{2}\equiv1\pmod p$ implies $(t-1)(t+1)\equiv0\pmod p$.

Pairing $j$ with $p-j$ for $1\le j\le m$, and using that $m$ is even,
we obtain
\[
 (p-1)!
 =\prod_{j=1}^{m}j(p-j)
 \equiv(-1)^{m}(m!)^{2}
 =(m!)^{2}\pmod p.
\]
Together with \eqref{eq:wilson}, this gives
\begin{equation}\label{eq:sqrt-minus-one}
 r^{2}\equiv-1\pmod p,
 \qquad r\equiv m!\pmod p.
\end{equation}

Let $M=\lfloor\sqrt p\rfloor$.  Since a prime is not a perfect square,
\[
 M^{2}<p<(M+1)^{2}.
\]
There are $(M+1)^{2}>p$ pairs $(u,v)$ with
$0\le u,v\le M$, but only $p$ residue classes of $u-rv$ modulo $p$.
By the pigeonhole principle, two distinct such pairs have the same
residue.  Subtracting them yields integers $A,B$, not both zero, such
that
\[
 |A|\le M,
 \qquad |B|\le M,
 \qquad A\equiv rB\pmod p.
\]
By \eqref{eq:sqrt-minus-one},
\[
 A^{2}+B^{2}\equiv(r^{2}+1)B^{2}\equiv0\pmod p.
\]
Moreover,
\[
 0<A^{2}+B^{2}\le2M^{2}<2p.
\]
Thus $A^{2}+B^{2}$ is a positive multiple of $p$ strictly smaller than
$2p$, so it equals $p$.
\end{proof}

\begin{theorem}[Two-square criterion]\label{thm:two-square}
For a positive integer $N$, the following conditions are equivalent:
\begin{enumerate}[label=\textup{(\roman*)},leftmargin=2.5em]
 \item $N$ is a sum of two rational squares;
 \item $N$ is a sum of two integer squares;
 \item $\ord_{q}(N)$ is even for every prime $q\equiv3\pmod4$.
\end{enumerate}
\end{theorem}

\begin{proof}
Condition (ii) immediately implies (i).  Suppose (i) holds, say
\[
 N=\left(\frac{A}{M}\right)^{2}
  +\left(\frac{B}{M}\right)^{2}
\]
for integers $A,B$ and a positive integer $M$.  Then
\[
 NM^{2}=A^{2}+B^{2}.
\]
For every prime $q\equiv3\pmod4$, Corollary~\ref{cor:valuation-even}
gives
\[
 \ord_{q}(N)+2\ord_{q}(M)
 =\ord_{q}(A^{2}+B^{2})\equiv0\pmod2.
\]
Hence $\ord_{q}(N)$ is even, proving (iii).

Now assume (iii).  In the prime factorisation of $N$, every prime
$q\equiv3\pmod4$ therefore occurs as $q^{2e}$, which is
$(q^{e})^{2}+0^{2}$.  Also $2=1^{2}+1^{2}$, and every prime
$p\equiv1\pmod4$ is a sum of two integer squares by
Lemma~\ref{lem:fermat-two-square}.  Finally, the identity
\begin{equation}\label{eq:brahmagupta}
 (A^{2}+B^{2})(C^{2}+D^{2})
 =(AC-BD)^{2}+(AD+BC)^{2}
\end{equation}
shows that a product of sums of two integer squares is again such a
sum.  Repeated application of \eqref{eq:brahmagupta} to the prime
factorisation of $N$ proves (ii).
\end{proof}

\section{A rational parametrisation of every circle fibre}

\begin{lemma}[The rational unit circle]\label{lem:unit-circle}
The map
\begin{equation}\label{eq:unit-param}
 \Phi:\Pone(\Q)\longrightarrow
 \{(A,B)\in\Q^{2}:A^{2}+B^{2}=1\},
 \qquad
 [u:v]\longmapsto
 \left(\frac{u^{2}-v^{2}}{u^{2}+v^{2}},
       \frac{2uv}{u^{2}+v^{2}}\right)
\end{equation}
is a bijection.
\end{lemma}

\begin{proof}
A nonzero rational pair $(u,v)$ cannot satisfy $u^{2}+v^{2}=0$, so the
map is well-defined.  It is invariant under replacing $(u,v)$ by
$(\lambda u,\lambda v)$ with $\lambda\ne0$, and direct expansion shows
that its image lies on the unit circle.

Conversely, let $(A,B)\in\Q^{2}$ satisfy $A^{2}+B^{2}=1$.  If $A=-1$,
then $B=0$, and $(A,B)=\Phi([0:1])$.  If $A\ne-1$, take
\[
 [u:v]=[1+A:B].
\]
Since
\[
 (1+A)^{2}+B^{2}=2(1+A),
\]
a direct calculation gives
\[
 \frac{(1+A)^{2}-B^{2}}{(1+A)^{2}+B^{2}}=A,
 \qquad
 \frac{2(1+A)B}{(1+A)^{2}+B^{2}}=B.
\]
This also gives an explicit inverse to $\Phi$, so the map is bijective.
\end{proof}

\begin{lemma}[A circle with a rational base point]\label{lem:general-circle}
Let $N>0$, and let $C,D\in\Q$ satisfy $C^{2}+D^{2}=N$.  Then the map
\begin{align}
 \Psi_{C,D}:\Pone(\Q)&\longrightarrow
 \{(R,S)\in\Q^{2}:R^{2}+S^{2}=N\},\label{eq:general-circle-map}\\
 [u:v]&\longmapsto
 \left(
 \frac{C(u^{2}-v^{2})-2Duv}{u^{2}+v^{2}},
 \frac{D(u^{2}-v^{2})+2Cuv}{u^{2}+v^{2}}
 \right)
\end{align}
is a bijection.
\end{lemma}

\begin{proof}
Let $(A,B)=\Phi([u:v])$ as in Lemma~\ref{lem:unit-circle}.  The displayed
map is equivalently
\begin{equation}\label{eq:rotation-matrix}
 \binom{R}{S}
 =
 \begin{pmatrix}C&-D\\D&C\end{pmatrix}
 \binom{A}{B}.
\end{equation}
Consequently
\[
 R^{2}+S^{2}=(C^{2}+D^{2})(A^{2}+B^{2})=N.
\]
Conversely, if $R^{2}+S^{2}=N$, define
\begin{equation}\label{eq:rotation-inverse}
 A=\frac{CR+DS}{N},
 \qquad
 B=\frac{CS-DR}{N}.
\end{equation}
Then $A,B\in\Q$ and
\[
 A^{2}+B^{2}
 =\frac{(C^{2}+D^{2})(R^{2}+S^{2})}{N^{2}}=1.
\]
Equations \eqref{eq:rotation-matrix} and
\eqref{eq:rotation-inverse} are inverse transformations.  The result
therefore follows from the bijectivity of $\Phi$.
\end{proof}

\section{Proof of the complete classification}

\begin{proposition}[Bijection with a two-square equation]\label{prop:fibre-bijection}
Let $a,b\in\Z_{>0}$ be coprime, let $\eps\in\{\pm1\}$, and put
$x=\eps a/b$.  The assignments
\begin{equation}\label{eq:RS-definition}
 R=b^{2}(2y-x^{2})=2b^{2}y-a^{2},
 \qquad
 S=2b^{2}z
\end{equation}
give a bijection
\[
 \mathcal S_{\eps a/b}(\Q)
 \longleftrightarrow
 \{(R,S)\in\Q^{2}:R^{2}+S^{2}=N_{a,b}\}.
\]
The inverse map is
\begin{equation}\label{eq:RS-inverse}
 y=\frac{a^{2}+R}{2b^{2}},
 \qquad
 z=\frac{S}{2b^{2}}.
\end{equation}
\end{proposition}

\begin{proof}
Multiplying \eqref{eq:norm-equation} by $b^{4}$ gives
\[
 \bigl(b^{2}(2y-x^{2})\bigr)^{2}
 +(2b^{2}z)^{2}
 =b^{4}(x^{4}-4)=a^{4}-4b^{4}.
\]
This is precisely $R^{2}+S^{2}=N_{a,b}$.  The formulas in
\eqref{eq:RS-definition} are visibly inverted by
\eqref{eq:RS-inverse}, proving the bijection.
\end{proof}

\begin{proof}[Proof of Theorem~\ref{thm:main}]
By Proposition~\ref{prop:fibre-bijection}, the fibre is nonempty if and
only if the positive integer $N_{a,b}$ is a sum of two rational squares.
The positivity is necessary because Proposition~\ref{prop:completion}
shows $x^{4}>4$.  Theorem~\ref{thm:two-square} proves that nonemptiness is
equivalent to \eqref{eq:criterion}, and also guarantees integers $C,D$
satisfying \eqref{eq:CD}.

For such $C,D$, Lemma~\ref{lem:general-circle} parametrises every pair
$(R,S)$ with $R^{2}+S^{2}=N_{a,b}$ exactly once.  Substitution into
\eqref{eq:RS-inverse} gives precisely
\eqref{eq:param-y}--\eqref{eq:param-z}.  Thus the asserted formulas are
complete and introduce no extraneous solutions.

It remains to prove the parity assertion \eqref{eq:parity-conclusion};
this is Proposition~\ref{prop:parity} below.  Once it is known, a
rational solution with integral $x$ is impossible, because the reduced
denominator of an integer is $1$, whereas every admissible denominator
is even.  In particular no triple in $\mathcal S(\Z)$ exists.
\end{proof}

\section{The numerator--denominator obstruction}

We use the elementary observation that every positive integer
congruent to $3$ modulo $4$ has at least one prime congruent to $3$
modulo $4$ occurring to odd exponent.  Indeed, if all such exponents
were even, its odd part would be congruent to $1$ modulo $4$.

\begin{proposition}\label{prop:parity}
Let $a,b\in\Z_{>0}$ be coprime.  If $N_{a,b}>0$ and every prime
$q\equiv3\pmod4$ occurs in $N_{a,b}$ with even exponent, then $a$ is odd
and $b$ is even.
\end{proposition}

\begin{proof}
Suppose first that $a$ is even.  Coprimality forces $b$ to be odd.  Write
$a=2k$.  Then
\begin{equation}\label{eq:a-even-factor}
 N_{a,b}=4(4k^{4}-b^{4}).
\end{equation}
The positive odd integer $4k^{4}-b^{4}$ is congruent to $3$ modulo $4$.
It therefore contains a prime $q\equiv3\pmod4$ to odd exponent.  Since
$q$ is odd, \eqref{eq:a-even-factor} shows that the same exponent occurs
in $N_{a,b}$, a contradiction.

It remains to exclude the possibility that both $a$ and $b$ are odd.
Since $N_{a,b}>0$, one has $a^{2}>2b^{2}$, and
\begin{equation}\label{eq:factor-N}
 N_{a,b}=(a^{2}-2b^{2})(a^{2}+2b^{2}).
\end{equation}
Both factors are positive and congruent to $3$ modulo $4$.  They are
coprime.  Indeed, any common divisor $d$ is odd and divides
\[
 (a^{2}+2b^{2})+(a^{2}-2b^{2})=2a^{2}
\]
and
\[
 (a^{2}+2b^{2})-(a^{2}-2b^{2})=4b^{2}.
\]
Thus $d\mid a^{2}$ and $d\mid b^{2}$, whence $d=1$ because
$\gcd(a,b)=1$.

Each factor in \eqref{eq:factor-N}, being $3$ modulo $4$, contains a
prime congruent to $3$ modulo $4$ to odd exponent.  Coprimality of the
factors prevents that exponent from being altered by the other factor.
Hence $N_{a,b}$ has a prime congruent to $3$ modulo $4$ to odd exponent,
again a contradiction.  Therefore $a$ is odd and $b$ is even.
\end{proof}

\begin{corollary}[Factorwise form of the criterion]\label{cor:factorwise}
Let $a,b\in\Z_{>0}$ be coprime.  The fibre above $x=\pm a/b$ is nonempty
over $\Q$ if and only if
\begin{enumerate}[label=\textup{(\alph*)},leftmargin=2.5em]
 \item $a$ is odd and $b$ is even;
 \item $a^{2}>2b^{2}$;
 \item each of $a^{2}-2b^{2}$ and $a^{2}+2b^{2}$ is a sum of two
 integer squares.
\end{enumerate}
\end{corollary}

\begin{proof}
Assume first that the fibre is nonempty.  Conditions (a) and (b) follow
from Proposition~\ref{prop:parity} and $N_{a,b}>0$.  Put
\[
 A=a^{2}-2b^{2},
 \qquad
 B=a^{2}+2b^{2}.
\]
Because $a$ is odd and $b$ is even, both $A$ and $B$ are odd.  The same
greatest-common-divisor argument used in Proposition~\ref{prop:parity}
shows $\gcd(A,B)=1$.  Since $AB=N_{a,b}$ and every prime
$q\equiv3\pmod4$ occurs in $AB$ with even exponent, coprimality implies
that it occurs with even exponent in each of $A$ and $B$.  Theorem
\ref{thm:two-square} proves (c).

Conversely, if (a)--(c) hold, then $N_{a,b}=AB>0$, and the product
identity \eqref{eq:brahmagupta} shows that $N_{a,b}$ is a sum of two
integer squares.  Proposition~\ref{prop:fibre-bijection} then gives a
rational point in the fibre.
\end{proof}

\begin{corollary}\label{cor:no-integral-x}
No rational solution has $x\in\Z$.  Consequently
\[
 \mathcal S(\Z)=\varnothing.
\]
\end{corollary}

\begin{proof}
If $x\in\Z$, its expression $x=\eps a/b$ in lowest terms has $b=1$.
Proposition~\ref{prop:parity} requires $b$ to be even, which is
impossible.
\end{proof}

\section{An independent congruence proof of the integral obstruction}

The preceding proof gives the stronger statement that integral $x$ is
impossible even when $y$ and $z$ are merely rational.  We record a
second proof of the weaker assertion $\mathcal S(\Z)=\varnothing$ that
uses only divisibility and congruences.

\begin{lemma}\label{lem:divisor-z2plus1}
Let $z\in\Z$, and let $d>0$ divide $z^{2}+1$.  Then
\begin{equation}\label{eq:d-form}
 d=2^{\delta}m,
 \qquad \delta\in\{0,1\},
\end{equation}
where every prime divisor of the odd integer $m$ is congruent to $1$
modulo $4$.  In particular, $m\equiv1\pmod4$.
\end{lemma}

\begin{proof}
Let $p$ be an odd prime divisor of $z^{2}+1$.  Then $p\nmid z$ and
$z^{2}\equiv-1\pmod p$.  If $p\equiv3\pmod4$, the argument in
Lemma~\ref{lem:3mod4-divides-both}, applied to $z^{2}+1^{2}$, would force
$p\mid1$, which is impossible.  Hence $p\equiv1\pmod4$.

For the power of $2$, if $z$ is even then $z^{2}+1$ is odd.  If $z$ is
odd, then $z^{2}\equiv1\pmod8$, so $z^{2}+1\equiv2\pmod8$ and its
$2$-adic valuation is exactly $1$.  Every positive divisor therefore
has the form \eqref{eq:d-form}.  Since all prime divisors of $m$ are
$1$ modulo $4$, so is $m$.
\end{proof}

\begin{theorem}[Direct integral obstruction]\label{thm:direct-integral}
There are no integers $x,y,z$ satisfying
\[
 y^{2}-x^{2}y+z^{2}+1=0.
\]
\end{theorem}

\begin{proof}
Assume an integral solution exists.  Equation \eqref{eq:product-form}
shows that $y>0$ and
\[
 y\mid z^{2}+1.
\]
By Lemma~\ref{lem:divisor-z2plus1}, write
$y=2^{\delta}m$, where $\delta\in\{0,1\}$ and $m\equiv1\pmod4$.

First suppose $\delta=0$.  Then $y$ is odd and $y\equiv1\pmod4$.
Rearrange the equation as
\begin{equation}\label{eq:integer-rearranged}
 x^{2}y=y^{2}+z^{2}+1.
\end{equation}
If $x$ were even, the left side would be $0$ modulo $4$, whereas the
right side would be $2$ or $3$ modulo $4$, according as $z$ is even or
odd.  Thus $x$ is odd.  If $z$ were even, the right side of
\eqref{eq:integer-rearranged} would be $2$ modulo $4$, while the left
side would be odd.  Hence $z$ is odd.  Reducing
\eqref{eq:integer-rearranged} modulo $4$ now gives
\[
 y\equiv1+1+1\equiv3\pmod4,
\]
contradicting $y\equiv1\pmod4$.

Now suppose $\delta=1$.  Then $y=2m\equiv2\pmod8$.  Since $y$ divides
$z^{2}+1$, the latter is even, so $z$ is odd.  Reducing the right side
of \eqref{eq:integer-rearranged} modulo $8$ gives
\[
 y^{2}+z^{2}+1\equiv4+1+1\equiv6\pmod8.
\]
If $x$ is even, then $x^{2}y\equiv0\pmod8$; if $x$ is odd, then
$x^{2}y\equiv2\pmod8$.  Neither possibility is $6$ modulo $8$, a final
contradiction.
\end{proof}

\section{Explicit infinite families}

\begin{example}[The complete fibre over $x=\pm3/2$]\label{ex:three-halves}
Take $a=3$ and $b=2$.  Then
\[
 N_{3,2}=3^{4}-4\cdot2^{4}=17=1^{2}+4^{2}.
\]
With $C=1$, $D=4$, and $[u:v]=[1:t]$, formulas
\eqref{eq:param-y}--\eqref{eq:param-z} become
\begin{align}
 x&=\pm\frac32,\label{eq:explicit-x}\\
 y&=\frac{4t^{2}-4t+5}{4(t^{2}+1)},\label{eq:explicit-y}\\
 z&=\frac{-2t^{2}+t+2}{4(t^{2}+1)},\label{eq:explicit-z}
\end{align}
where $t\in\Q$.  The missing projective parameter $[0:1]$, customarily
written $t=\infty$, gives
\[
 (x,y,z)=\left(\pm\frac32,1,-\frac12\right).
\]
By Theorem~\ref{thm:main}, these formulas give not merely infinitely
many examples but every rational point in either fibre $x=3/2$ or
$x=-3/2$.  For instance, $t=0$ gives
\[
 \left(\frac32,\frac54,\frac12\right).
\]
\end{example}

The preceding example already proves that $\mathcal S(\Q)$ is infinite.
The next construction shows that infinitely many distinct rational
$x$-coordinates occur.

\begin{proposition}[A Pell subfamily]\label{prop:pell-family}
Define positive integer pairs $(a_n,b_n)$ recursively by
\begin{equation}\label{eq:pell-recurrence}
 (a_1,b_1)=(3,2),
 \qquad
 a_{n+1}=3a_n+4b_n,
 \qquad
 b_{n+1}=2a_n+3b_n.
\end{equation}
Then, for every $n\ge1$,
\begin{equation}\label{eq:pell-invariant}
 a_n^{2}-2b_n^{2}=1,
\end{equation}
$\gcd(a_n,b_n)=1$, and the fibre above $x=\pm a_n/b_n$ is nonempty.
More explicitly, for every $[u:v]\in\Pone(\Q)$, with
$U=u^{2}+v^{2}$,
\begin{align}
 x&=\pm\frac{a_n}{b_n},\label{eq:pell-x}\\
 y&=\frac{a_n^{2}U+(u^{2}-v^{2})-4b_nuv}
 {2b_n^{2}U},\label{eq:pell-y}\\
 z&=\frac{2b_n(u^{2}-v^{2})+2uv}
 {2b_n^{2}U}\label{eq:pell-z}
\end{align}
solves the equation.  The positive numbers $a_n/b_n$ are pairwise
distinct.
\end{proposition}

\begin{proof}
A direct calculation shows that the recurrence preserves the quadratic
form $a^{2}-2b^{2}$:
\begin{align*}
 &(3a+4b)^{2}-2(2a+3b)^{2}\\
 &\qquad=9a^{2}+24ab+16b^{2}
 -8a^{2}-24ab-18b^{2}\\
 &\qquad=a^{2}-2b^{2}.
\end{align*}
Since $3^{2}-2\cdot2^{2}=1$, induction proves
\eqref{eq:pell-invariant}.  Any common divisor of $a_n$ and $b_n$ would
have square dividing the left side of \eqref{eq:pell-invariant}, so it
must be $1$.  The recurrence also gives
$b_{n+1}=2a_n+3b_n>b_n$, so the sequence $(b_n)$ is strictly increasing.

From \eqref{eq:pell-invariant},
\begin{align}
 N_{a_n,b_n}
 &=(a_n^{2}-2b_n^{2})(a_n^{2}+2b_n^{2})\notag\\
 &=1\cdot(1+4b_n^{2})
 =1^{2}+(2b_n)^{2}.\label{eq:pell-N}
\end{align}
Thus one may choose $C=1$ and $D=2b_n$ in
Theorem~\ref{thm:main}; formulas
\eqref{eq:pell-x}--\eqref{eq:pell-z} are exactly the resulting
parametrisation.

Finally,
\[
 \left(\frac{a_n}{b_n}\right)^{2}
 =2+\frac{1}{b_n^{2}}.
\]
Since the positive integers $b_n$ strictly increase, these squares, and
hence the positive ratios $a_n/b_n$, are pairwise distinct.
\end{proof}

The first few pairs are
\[
 (a_n,b_n)=(3,2),\ (17,12),\ (99,70),\ (577,408),\ldots.
\]
Thus rational points occur on infinitely many distinct circle fibres,
while every such $x$ necessarily has an even reduced denominator.

\end{document}
